\section{Conclusion}
\label{sec:conclusion}

This paper presented an empirical investigation into lesion-grounded multi-task deep learning for rice leaf disease recognition using an EfficientNet-B0 backbone. By combining six-class pathology classification with bounding-box localization supervision and post-processing calibration, the proposed Phase 3B model achieves 88.96\% accuracy, 0.8751 Macro F1, a Localization F1 of 0.0905, and a Mean Matched IoU of 0.6423 on the primary RiceLeafDiseaseBD benchmark. Compared to the uncalibrated prototype, refinement reduced healthy leaf false positive detections from 21.52\% to 10.97\% and decreased background border attention from 18.5\% to 9.2\%. Quantitative XAI validation against 4,348 independent RiceSeg5932 masks confirmed directional Pointing Game gains on correctly classified samples alongside reduced attribution entropy. The findings establish that weak bounding-box supervision successfully imparts spatial lesion localization capability and reduces background shortcut reliance while highlighting a measurable, limited classification trade-off relative to pure classification baselines.

Future work includes exploring multi-scale feature pyramids (FPN) for dense lesion resolution, investigating temperature-scaled uncertainty calibration for selective rejection, and conducting prospective field evaluations across diverse agro-ecological zones.
